% DO NOT MODIFY THIS FIILE
% ----------------------------------------------
\usepackage{relsize}
\usepackage{tabulary}
\usepackage[dvipsnames, table]{xcolor}
\usepackage[a4paper, portrait, margin=2cm, top=1.8cm, left=1.8cm, bottom=2cm]{geometry}
\usepackage{titlesec}
\usepackage{tikz}
    \usetikzlibrary{shapes.misc}
    \usetikzlibrary{calc}
\usepackage{setspace}
\usepackage{fontspec}
\usepackage[most]{tcolorbox}
\usepackage{varwidth}
\usepackage[inline]{enumitem}
\usepackage{multicol}
\usepackage{fancyhdr}
\usepackage{booktabs, multirow, tabularx}
\usepackage{mathtools}
\usepackage{tasks}
\usepackage{anyfontsize}
\usepackage[export]{adjustbox}          % http://ctan.org/pkg/adjustbox
\usepackage{bookmark}



\newcommand{\greentheme}{
    \definecolor{primarythemecol}{RGB}{39, 165, 85}
}
\newcommand{\purpureustheme}{
    \definecolor{primarythemecol}{RGB}{157, 68, 181}
}
\newcommand{\saffrontheme}{
    \definecolor{primarythemecol}{RGB}{220, 177, 4}
}
\newcommand{\bluetheme}{
    \definecolor{primarythemecol}{RGB}{38,106,234}
}




% Multi-column enumeration
\newcounter{taskscounter}
\settasks{
    style=enumerate,
    counter=taskscounter,
    label-width={22pt},
    item-indent={15pt},
    label-align=right,
    label=\textbf{\alph*},
    before-skip = -\parskip ,           % undo paragraph skip
    after-skip = -\parskip ,            % undo paragraph skip
    after-item-skip = -\parskip+1mm     % undo paragraph skip
%   debug=true                          % useful for fine-tuning or debugging
}




% Define dynamic titlebars for sections/exercises
% Heavily modified from a post on StackExchange
% https://tex.stackexchange.com/questions/34288/how-to-place-a-shaded-box-around-a-section-label-and-name
\newlength{\sectionnumberwidth}
\newcommand\titlebar{\hspace*{-0.1cm}
    \tikz[baseline, trim right=3.1cm] {
        \settowidth{\sectionnumberwidth}{
            \pgfinterruptpicture
                \textbf{\sffamily\thesection.\thesubsection}
            \endpgfinterruptpicture
        }
        \fill [primarythemecol!30!white, rounded corners=0.15cm] (2.5cm,-1.28ex) rectangle (\textwidth-\sectionnumberwidth+2.65cm,2.615ex);
        \node [
            fill=primarythemecol,
            text=white,
            anchor= base east,
            rounded rectangle,
            inner sep = 2mm,
            minimum height=3.89ex] at (3cm,0) {
            {\bfseries\thesection.\thesubsection}
        };
    }
}
\newcommand\exercisebar{\hspace*{-0.1cm}
    \tikz[baseline, trim right=3.1cm] {
        \settowidth{\sectionnumberwidth}{
            \pgfinterruptpicture
                \bfseries{\sffamily\thesection\thesubsubsection}
            \endpgfinterruptpicture
        }
        \fill [red!25, rounded corners=0.15cm] (2.5cm,-1.22ex) rectangle (\textwidth-\sectionnumberwidth+2.65cm,2.67ex);
        \node [
            fill=cherryred!80!white,
            text=white,
            anchor= south east,
            rounded rectangle,
            inner sep = 2mm,
            minimum height=3.9ex] at (2.6cm,-2.62mm) {
            {\bfseries\thesection\thesubsubsection}
        };
    }
}

% Implement the custom titlebars
\titleformat{\section}{\large\sffamily}{\titlebar}{0.1cm}{}
\titleformat{\subsection}{\large\sffamily}{\titlebar}{0.1cm}{}
\titleformat{\subsubsection}{\large\sffamily}{\exercisebar}{0.1cm}{}

% Define phantom subsections/subsubsections
\newcommand{\fakesubsection}[1]{%
  \par\refstepcounter{subsection}    % Increase subsection counter
  \subsectionmark{#1}                % Add subsection mark (header)
  \addcontentsline{toc}{subsection}{\protect\numberline{\thesubsection}#1}% Add subsection to ToC
  % Add more content here, if needed.
}
\newcommand{\fakesubsubsection}[1]{%
  \par\refstepcounter{subsubsection}    % Increase subsection counter
  \subsectionmark{#1}                   % Add subsection mark (header)
  \addcontentsline{toc}{subsubsection}{\protect\numberline{\thesubsubsection}#1}% Add subsection to ToC
  % Add more content here, if needed.
}




% Modified from Section 10.2.3 of tcolorbox manual, found here:
% https://mirror.ox.ac.uk/sites/ctan.org/macros/latex/contrib/tcolorbox/tcolorbox.pdf
\newtcolorbox{mybox2}[2][]{
    enhanced,
    breakable,
    before skip=2mm,after skip=2mm,
    colback=black!5,
    colframe=black!50,
    boxrule=0.2mm,
    attach boxed title to top left={
        xshift=1cm,yshift*=1mm-\tcboxedtitleheight},
        varwidth boxed title*=-3cm,
        boxed title style={frame code={
            \path[fill=tcbcolback!30!black]
              ([yshift=-1mm,xshift=-1mm]frame.north west)
                arc[start angle=0,end angle=180,radius=1mm]
              ([yshift=-1mm,xshift=1mm]frame.north east)
                arc[start angle=180,end angle=0,radius=1mm];
            \path[left color=tcbcolback!60!black,right color=tcbcolback!60!black,
              middle color=tcbcolback!80!black]
              ([xshift=-2mm]frame.north west) -- ([xshift=2mm]frame.north east)
              [rounded corners=1mm]-- ([xshift=1mm,yshift=-1mm]frame.north east)
              -- (frame.south east) -- (frame.south west)
              -- ([xshift=-1mm,yshift=-1mm]frame.north west)
              [sharp corners]-- cycle;
            },interior engine=empty,
          },
    fonttitle=\sffamily\bfseries,
    title={#2},
    before upper = \sffamily,
    #1
}




\tcbset{%
    sidebarwithnum/.style={%
        enhanced,
        breakable,
        rounded corners,
        toprule=0pt, rightrule=0pt, bottomrule=0pt, leftrule=1mm,
        colback=#1!5, colframe=#1!90!black, coltitle=#1!90!black,
        fonttitle=\large\sffamily\sbseries,
        detach title,
        before upper={\tcbtitle\quad},
        fontupper=\linespread{1.2}\selectfont
    },
    sidebarwithoutnum/.style={%
        enhanced,
        breakable,
        separator sign none,
        rounded corners,
        toprule=0pt, rightrule=0pt, bottomrule=0pt, leftrule=1mm,
        colback=#1!5, colframe=#1!80!black, coltitle=#1!80!black,
        fonttitle=\large\sffamily\sbseries,
        detach title,
        before upper={\tcbtitle\quad\sffamily},
        fontupper=\linespread{1.2}\selectfont
    },
}
\newtcbtheorem[auto counter]{examplebox}{Example}
{sidebarwithnum=primarythemecol}{ex}
\newtcbtheorem{practice}{Independent Practice}
{sidebarwithnum=practiceorange}{pr}
\newtcbtheorem{note}{#1\hspace{-1mm}}
{sidebarwithoutnum=notecolor}{nt}




% Custom Colour Definitions
\definecolor{practiceorange}{RGB}{252, 191, 0}
\definecolor{notecolor}{RGB}{0, 0, 0}
\definecolor{cherryred}{RGB}{191,0,0}

\definecolor{ehclogoblue}{RGB}{0,191,214}
\definecolor{ehclogogray}{RGB}{34, 34, 34}
\definecolor{ehclogoorange}{RGB}{255, 107, 0}
\definecolor{separatorgray}{RGB}{210, 210, 210}




% Changing various spacings
% ----------------------------------------------
\setstretch{1.1}                                            % Globally adjust bullet spacing
\setlength{\parindent}{0cm}                                 % Removes paragraph indent
\setlength{\multicolsep}{6.0pt plus 2.0pt minus 1.5pt}      % Halves whitespace before multicol environment
\setlength\tabcolsep{0pt}                                   % Removes default space between columns in table

% Change various spacings of bulleted/numbered lists
\setlist[itemize]{topsep=2pt, itemsep=-2pt}
\setlist[enumerate]{topsep=-2pt, partopsep=4mm, itemsep=0pt}

% Change spacing of aligned mathematics
\AtBeginDocument{
    \setlength\abovedisplayskip{3pt}
    %\setlength{\belowdisplayskip}{0pt}
    \setlength\belowdisplayskip{7pt}
    %\setlength{\belowdisplayshortskip}{0pt}
}

\AtEndEnvironment{tabularx}{\vspace{-4mm}}

\newcommand{\Sum} [2] {\the\numexpr #1 + #2 \relax \\}
\AtBeginEnvironment{examplebox}{
    \fakesubsubsection{Example \Sum{\number\value{tcb@cnt@examplebox}}{1} }
}

% Function for adding vspace above underbrace
\newcommand*\addunderbracespace[1]{\vrule width0pt height0pt depth#1\relax}

% Define custom struts for aligning square
% #1 the height, #2 the depth
% Credit: https://tex.stackexchange.com/questions/14602/set-strut-height
\def\mystrut(#1,#2){\vrule height #1pt depth #2pt width 0pt}




% Custom textbook font definitions
\setsansfont{AptiferSansPro}[
    Path = fonts/,
    Extension = .ttf,
    UprightFont = *-Regular,
    BoldFont = *-Semibold,
    BoldItalicFont = *-SemiboldItalic,
    FontFace = {l}{n}{*-Light},
    FontFace = {l}{it}{*-LightItalic},
    FontFace = {sb}{n}{*-Medium},
    FontFace = {sb}{it}{*-MediumItalic},
    FontFace = {eb}{n}{*-Black},
    FontFace = {eb}{it}{*-BlackItalic},
]
\DeclareRobustCommand{\sbseries}{\fontseries{sb}\selectfont}
\DeclareTextFontCommand{\textsb}{\sbseries}
\DeclareRobustCommand{\ebseries}{\fontseries{eb}\selectfont}
\DeclareTextFontCommand{\texteb}{\ebseries}
\DeclareRobustCommand{\ltseries}{\fontseries{l}\selectfont}
\DeclareTextFontCommand{\textlt}{\ltseries}







% Define the easier get functions for sections and subsections
\renewcommand*{\thesection}{\arabic{section}}
\renewcommand*{\thesubsection}{\arabic{subsection}}
\renewcommand*{\thesubsubsection}{\Alph{subsubsection}}
\newcommand\getcurrentref[1]{%
 \ifnumequal{\value{#1}}{0}
  {??}
  {\the\value{#1}}%
}
%\newcommand{\exercise}{\subsubsection}
\newcommand{\exercise}[1]{\subsubsection{Exercise \thesection\thesubsubsection {\footnotesize {\hspace{4mm}}#1}}}
\newcommand{\lesson}[1]{\subsection{#1}}


% Custom font size 'YUGE', larger than huge
\newcommand\YUGE{\fontsize{35}{35}\selectfont}




% Custom Footer Definitions
% ----------------------------------------------
\fancypagestyle{branded}{
    \fancyhf{} % clear all header and footer fields
    \fancyfoot[OC]{ %right hand side (odd pages)
        \vspace{-0.7em}
        \hfill {\sffamily\raggedright {\color{separatorgray}\getcurrentref{page}} \\}\vspace{-1.6em}
        \begin{minipage}{.3\linewidth}
            \color{separatorgray}\rule{\linewidth}{0.4pt}
        \end{minipage}
        \begin{minipage}{.1\linewidth}
            \centering
            \includegraphics[scale=0.4]{images/branding/brand-separator-icon}
        \end{minipage}
        \begin{minipage}{.3\linewidth}
            \color{separatorgray}\rule{\linewidth}{0.4pt}
        \end{minipage}
    }
    \fancyfoot[EC]{ %left hand side (even pages)
        \vspace{-0.7em}
        {\sffamily\raggedright {\color{separatorgray}\getcurrentref{page}} \\}\vspace{-1.6em}
        \begin{minipage}{.3\linewidth}
            \color{separatorgray}\rule{\linewidth}{0.4pt}
        \end{minipage}
        \begin{minipage}{.1\linewidth}
            \centering
            \includegraphics[scale=0.4]{images/branding/brand-separator-icon}
        \end{minipage}
        \begin{minipage}{.3\linewidth}
            \color{separatorgray}\rule{\linewidth}{0.4pt}
        \end{minipage}
    }

    \renewcommand{\headrulewidth}{0pt}
    \renewcommand{\footrulewidth}{0pt}
}
\fancypagestyle{plain}{
    \fancyhf{} % clear all header and footer fields
    \fancyfoot[OC]{ %right hand side
        \hfill\sffamily\raggedright\getcurrentref{page}
    }
    \fancyfoot[EC]{ %left hand side
        \sffamily\raggedright \getcurrentref{page}
    }

    \renewcommand{\headrulewidth}{0pt}
    \renewcommand{\footrulewidth}{0pt}
}
\fancypagestyle{attribution}{
    \fancyhf{} % clear all header and footer fields
    \fancyfoot[L]{ %left hand side
        \vspace{-3.0em}\color{separatorgray}\small\sffamily
        \bookletsubtitle \\
        {\rmfamily{\LaTeX}} Template created by: Ellis Dickinson \\
        Document compiled on: \today
    }
    \renewcommand{\headrulewidth}{0pt}
    \renewcommand{\footrulewidth}{0pt}
}



% Metadata Processing
% ----------------------------------------------
\newcommand{\formattedunittitle}{Unit \unitprefix\getcurrentref{section}}



% Additional tabular column types
% ----------------------------------------------
\newcolumntype{P}[1]{>{\centering\arraybackslash}p{#1}}
\newcolumntype{C}{>{\centering\arraybackslash}X}



% Additional Tabular Lines
% ----------------------------------------------
\makeatletter
\def\thickhline{%
  \noalign{\ifnum0=`}\fi\hrule \@height \thickarrayrulewidth \futurelet
   \reserved@a\@xthickhline}
\def\@xthickhline{\ifx\reserved@a\thickhline
               \vskip\doublerulesep
               \vskip-\thickarrayrulewidth
             \fi
      \ifnum0=`{\fi}}
\makeatother

\newlength{\thickarrayrulewidth}
\setlength{\thickarrayrulewidth}{2\arrayrulewidth}
